\documentclass[11pt,a4paper]{article}

\usepackage[utf8]{inputenc}
\usepackage[T1]{fontenc}
\usepackage{lmodern}
\usepackage{geometry}
\usepackage{booktabs}
\usepackage{longtable}
\usepackage{array}
\usepackage{amsmath}
\usepackage{xcolor}
\usepackage{hyperref}
\usepackage{enumitem}
\usepackage{tabularx}
\usepackage{makecell}

\geometry{margin=1in}
\hypersetup{
  colorlinks=true,
  linkcolor=blue,
  urlcolor=blue,
  citecolor=blue
}

\setlist[itemize]{topsep=2pt,itemsep=2pt,parsep=0pt}
\setlist[enumerate]{topsep=2pt,itemsep=2pt,parsep=0pt}
\renewcommand{\arraystretch}{1.15}

\newcommand{\code}[1]{\texttt{#1}}
\newcommand{\pathcode}[1]{{\small\nolinkurl{#1}}}
\newcommand{\metric}[1]{\textbf{#1}}

\title{BOAMP Grand Ouest Renewal-Linkage Project\\
\large Historical Session Documentation, Metrics, and Proposed Next Pipeline}
\author{Prepared for the BOAMP Data Science Internship Project}
\date{11 August 2026}

\begin{document}
\maketitle

\section*{Document Status}

\noindent\fbox{\parbox{0.94\textwidth}{%
\textbf{Superseded historical report.} This document records the project state on
11 August 2026, before the current v2 standardisation, episode reconstruction,
full-cohort linkage, Fellegi--Sunter comparison, expiry-aware sensitivity arm, and
survival analysis were completed. Statements below using ``current'', ``not yet'',
or ``next'' refer to that historical session and must not be read as the present
repository state. The current methodology is documented in
\pathcode{reports/boamp_methodology_chapter.tex}; the expiry extension is documented
in \pathcode{reports/expiry_aware_successor_selection_v1.tex}.}}

The raw EDA in this historical branch counts literal \code{nomacheteur} strings.
Consequently \code{Ville de Paris} and \code{VILLE DE PARIS} appear as different
buyers in its chart. This is a property of the raw source and the exact-string EDA,
not of the current v2 identity layer, which case-folds both to
\code{ville de paris}. Longer organisational suffixes still require SIREN or other
entity-resolution evidence.

\section*{Technical Summary}

This document consolidates a historical working session on the BOAMP public procurement project. The project objective is to study renewal behavior of digital public procurement contracts in Grand Ouest using BOAMP data. BOAMP does not provide an explicit ground-truth field stating that one contract renews another, so the project requires a record-linkage stage before survival analysis.

The recommended methodological position is conservative:

\begin{quote}
Use BOAMP to identify high-confidence, observable renewal links. Do not claim to recover every true renewal in the procurement market.
\end{quote}

The current project status is:

\begin{center}
\small
\begin{tabularx}{\textwidth}{p{0.30\textwidth}p{0.18\textwidth}X}
\toprule
Phase & Status & Main Artifact \\
\midrule
Raw BOAMP acquisition & Done & \pathcode{data/raw/boamp/} \\
Grand Ouest filtering & Done & \pathcode{boamp_2015_2025_grand_ouest_raw.csv} \\
EDA & Done & \pathcode{notebooks/boamp_full_raw_eda.ipynb} \\
Preprocessing / engineering & Done & Notebook 03 \\
Manual reference benchmark & Done & Notebook 04 \\
Candidate pair generation & Done for benchmark anchors & Notebook 05 \\
Baseline linkage comparison & Done & Notebook 06 \\
High-precision iteration & Done & Notebook 08 \\
Parameter diagnostics & Done & Notebook 09 \\
Full-corpus episode linkage & Not yet & Next implementation \\
Survival analysis & Not yet & After linkage is frozen \\
\bottomrule
\end{tabularx}
\end{center}

The current best linkage method is a transparent high-precision gated weighted score. On the locked test benchmark it improves precision and false-positive behavior compared with the first weighted baseline, but it is still not accurate enough to describe as a complete renewal detector.

\section*{Project Objective and Clean Flow}

The real internship objective is not simply to build many linkage algorithms. The project goal is:

\begin{quote}
Study renewal behavior and future procurement needs for digital public contracts in Grand Ouest using BOAMP data.
\end{quote}

The clean end-to-end flow is:

\begin{enumerate}
  \item Acquire raw BOAMP notices for 2015--2025.
  \item Restrict the analysis scope to Grand Ouest.
  \item Perform exploratory data analysis.
  \item Engineer buyer, date, text, CPV, amount, duration, and quality features.
  \item Prepare a manually validated benchmark because BOAMP has no explicit renewal label.
  \item Generate candidate successor pairs.
  \item Compare simple linkage methods.
  \item Tune parameters on \code{PILOT\_DEVELOPMENT}.
  \item Evaluate once on \code{LOCKED\_TEST}.
  \item Freeze one conservative linkage rule.
  \item Apply the frozen rule to the full Grand Ouest digital-contract population.
  \item Build renewal episodes and survival-analysis data.
  \item Run Kaplan-Meier, Cox, and sensitivity analyses.
\end{enumerate}

\section*{External Methodological Context}

The method follows standard record-linkage logic:

\begin{itemize}
  \item The Fellegi-Sunter tradition treats linkage as cumulative evidence across several fields rather than a single exact identifier.\footnote{\url{https://moj-analytical-services.github.io/splink/topic_guides/theory/fellegi_sunter.html}}
  \item Deterministic linkage is cheap and can achieve high precision, but usually has lower recall; probabilistic linkage allows the analyst to choose a threshold depending on the cost of false positives versus false negatives.\footnote{\url{https://moj-analytical-services.github.io/splink/topic_guides/theory/probabilistic_vs_deterministic.html}}
  \item Blocking or indexing is necessary in large record linkage because comparing all possible pairs is computationally infeasible.\footnote{\url{https://openresearch-repository.anu.edu.au/items/29ab1c2e-6b1f-4d61-8691-185fed831565}}
  \item Active learning and clerical review are useful when ground truth is limited; a small number of carefully chosen manual labels can improve linkage performance.\footnote{\url{https://pmc.ncbi.nlm.nih.gov/articles/PMC10072450/}}
  \item Survival analysis is appropriate once a renewal event and time-to-event variable are defined. Right-censored observations should be kept, not dropped.\footnote{\url{https://pmc.ncbi.nlm.nih.gov/articles/PMC10357905/}}
\end{itemize}

BOAMP itself is the official source used for the raw notices.\footnote{\url{https://boamp-datadila.opendatasoft.com/explore/dataset/boamp/api/}} CPV is an official hierarchical public procurement classification system, which justifies using CPV prefixes as evidence but not as the sole matching rule.\footnote{\url{https://single-market-economy.ec.europa.eu/single-market/public-procurement/digital-procurement/common-procurement-vocabulary_en}}

\section*{Data Scope and Raw Acquisition}

The main historical corpus covers:

\[
2015\text{-}01\text{-}01 \leq \texttt{dateparution} < 2026\text{-}01\text{-}01
\]

The Grand Ouest working scope is:

\begin{itemize}
  \item Bretagne
  \item Normandie
  \item Pays de la Loire
\end{itemize}

The Grand Ouest raw CSV summary is:

\begin{center}
\begin{tabular}{lr}
\toprule
Metric & Value \\
\midrule
Full BOAMP rows scanned & 1,620,712 \\
Grand Ouest rows & 226,314 \\
Raw fields preserved & 41 \\
Unique \code{idweb} & 226,314 \\
Duplicate \code{idweb} & 0 \\
Missing \code{idweb} & 0 \\
Minimum \code{dateparution} & 2015-03-02 \\
Maximum \code{dateparution} & 2025-12-31 \\
Records outside 2015--2025 range & 0 \\
\bottomrule
\end{tabular}
\end{center}

Rows by Grand Ouest region:

\begin{center}
\begin{tabular}{lr}
\toprule
Region & Notices \\
\midrule
Bretagne & 64,208 \\
Normandie & 86,180 \\
Pays de la Loire & 75,926 \\
\bottomrule
\end{tabular}
\end{center}

Rows by year:

\begin{center}
\begin{tabular}{rrrr}
\toprule
Year & Notices & Year & Notices \\
\midrule
2015 & 18,462 & 2021 & 20,827 \\
2016 & 20,349 & 2022 & 21,349 \\
2017 & 20,416 & 2023 & 22,012 \\
2018 & 20,865 & 2024 & 20,507 \\
2019 & 21,601 & 2025 & 20,754 \\
2020 & 19,172 & Total & 226,314 \\
\bottomrule
\end{tabular}
\end{center}

\section*{Preprocessing and Engineered Dataset}

The preprocessing notebook creates:

\begin{quote}
\pathcode{data/processed/boamp_grand_ouest/notices_engineered.parquet}
\end{quote}

It engineers identity, date, geography, buyer, text, CPV/theme, amount, duration, and quality features. Missing values are not blindly imputed. They are preserved and represented with explicit missingness flags.

\begin{center}
\begin{tabular}{lr}
\toprule
Metric & Value \\
\midrule
Engineered rows & 226,314 \\
Engineered fields & 74 \\
Unique \code{idweb} & 226,314 \\
Missing \code{idweb} & 0 \\
Duplicate \code{idweb} & 0 \\
Invalid \code{dateparution} & 0 \\
Outside historical range & 0 \\
Missing buyer key & 10 \\
Missing CPV & 39,993 \\
Digital candidate rows & 24,869 \\
Usable for linkage rows & 226,299 \\
\bottomrule
\end{tabular}
\end{center}

Digital candidate rows by segment:

\begin{center}
\begin{tabular}{lr}
\toprule
Segment & Rows \\
\midrule
Telecom / network & 8,655 \\
IT services & 7,337 \\
Software & 3,199 \\
Cybersecurity & 2,841 \\
Hardware / infrastructure & 1,479 \\
Cloud / hosting & 867 \\
Data / AI & 491 \\
\bottomrule
\end{tabular}
\end{center}

\section*{Manual Reference Benchmark}

Because BOAMP does not include a renewal label, the reference benchmark acts like a train/test split in normal machine learning:

\begin{itemize}
  \item \code{PILOT\_DEVELOPMENT}: used to choose thresholds, gates, and weight presets.
  \item \code{LOCKED\_TEST}: held out and used only once for final evaluation.
\end{itemize}

The current benchmark files are:

\begin{itemize}
  \item \pathcode{data/reference/BOAMP_Internship_Reference_120.csv}
  \item \pathcode{data/reference/BOAMP_Internship_Evaluation_Subset.csv}
  \item \pathcode{data/reference/BOAMP_Internship_Confirmed_Successor_Links.csv}
\end{itemize}

Benchmark summary:

\begin{center}
\begin{tabular}{lr}
\toprule
Metric & Value \\
\midrule
Reference anchors & 120 \\
Evaluation subset anchors & 94 \\
Confirmed successor link rows & 28 \\
\code{PILOT\_DEVELOPMENT} anchors & 16 \\
\code{LOCKED\_TEST} anchors & 78 \\
Anchor notice IDs missing in BOAMP & 0 \\
Successor notice IDs missing in BOAMP & 0 \\
\bottomrule
\end{tabular}
\end{center}

Evaluation outcome distribution:

\begin{center}
\begin{tabular}{lr}
\toprule
Outcome & Count \\
\midrule
\code{NO\_OBSERVED\_SUCCESSOR\_IN\_SCOPE} & 70 \\
\code{OBSERVED\_SUCCESSOR} & 23 \\
\code{MULTIPLE\_SUCCESSORS} & 1 \\
\bottomrule
\end{tabular}
\end{center}

\section*{Candidate Pair Generation}

Candidate generation is currently benchmark-focused. It creates possible future successors for reference anchors:

\begin{quote}
\pathcode{data/processed/boamp_grand_ouest/renewal_candidate_pairs.parquet}
\end{quote}

Candidate rules:

\begin{itemize}
  \item candidate date must be after anchor date;
  \item time gap between roughly 3 months and 8 years;
  \item same buyer key or similar normalized buyer name;
  \item same CPV/theme or strong text similarity;
  \item Grand Ouest only.
\end{itemize}

Candidate generation metrics:

\begin{center}
\begin{tabular}{lr}
\toprule
Metric & Value \\
\midrule
Candidate pairs & 8,095 \\
Reference anchors covered & 120 / 120 \\
Eligible evaluation anchors covered & 94 / 94 \\
\bottomrule
\end{tabular}
\end{center}

\section*{Linkage Methods Tested}

The following methods were implemented and compared:

\small
\begin{longtable}{p{0.34\textwidth}p{0.56\textwidth}}
\toprule
Method & Explanation \\
\midrule
\endhead
\makecell[l]{M0: no-link\\baseline} & Always predicts no successor. This is a sanity baseline. \\
\makecell[l]{M1: exact-rule\\baseline} & Accepts strict deterministic matches based on buyer and timing. \\
\makecell[l]{M2: nearest valid\\successor} & Chooses the nearest future valid candidate within the allowed time window. \\
\makecell[l]{M3: weighted\\rule score} & First main baseline. Uses weighted buyer, text, CPV, time, and geography evidence, then applies a threshold. \\
\makecell[l]{M4: word TF-IDF\\similarity} & Ranks candidates mainly by word-level TF-IDF similarity. \\
\makecell[l]{M5: char n-gram\\TF-IDF similarity} & Ranks candidates mainly by character n-gram TF-IDF similarity. \\
\makecell[l]{M6: high-precision\\gated score} & Current best direction. Uses a weighted score plus explicit acceptance gates to reduce false positives. \\
\bottomrule
\end{longtable}
\normalsize

\section*{Linkage Metrics}

The evaluation uses the following metrics:

\small
\begin{longtable}{p{0.34\textwidth}p{0.56\textwidth}}
\toprule
Metric & Meaning \\
\midrule
\endhead
\metric{Coverage rate / linking rate} & Share of anchors for which the model predicts at least one successor. \\
\metric{Precision@1} & Among anchors with a predicted top successor, share where the top prediction is correct. This is the most important automatic-link quality metric. \\
\metric{Recall@1} & Among anchors that truly have a successor in the benchmark, share where the true successor is found as top-1. \\
\metric{Recall@5} & Among anchors that truly have a successor, share where the true successor appears anywhere in the top 5 candidates. Useful for manual review. \\
\metric{Mean reciprocal rank} & Rewards ranking the true successor near the top of the candidate list. \\
\metric{False-positive rate on no-successor anchors} & Among anchors manually labelled as no observed successor, share where the model wrongly predicts a successor. This matters strongly for survival analysis. \\
\metric{No-link accuracy} & Among no-successor anchors, share where the model correctly predicts no link. \\
\metric{Predicted count} & Number of anchors where the model emits an automatic link. \\
\bottomrule
\end{longtable}
\normalsize

\section*{Baseline Versus High-Precision Results}

The first useful baseline was \code{M3\_weighted\_rule\_score}. It was too permissive: it linked more anchors, but many predicted renewals were false positives.

\begin{center}
\small
\begin{tabularx}{\textwidth}{Xrr}
\toprule
Metric on \code{LOCKED\_TEST} & M3 baseline & M6 high precision \\
\midrule
Anchors & 78 & 78 \\
Positive successor anchors & 19 & 19 \\
No-successor anchors & 59 & 59 \\
Coverage rate & 0.564 & 0.410 \\
Precision@1 & 0.250 & 0.438 \\
Recall@1 & 0.579 & 0.737 \\
Recall@5 & 0.684 & 0.842 \\
Mean reciprocal rank & 0.152 & 0.190 \\
False-positive rate & 0.441 & 0.237 \\
No-link accuracy & 0.559 & 0.763 \\
Predicted count & 44 & 32 \\
\bottomrule
\end{tabularx}
\end{center}

Interpretation:

\begin{itemize}
  \item The high-precision method improves \metric{precision@1} from 0.250 to 0.438.
  \item It reduces \metric{false-positive rate} from 0.441 to 0.237.
  \item It lowers \metric{coverage} from 0.564 to 0.410.
  \item This is the expected tradeoff: fewer links, but more reliable links.
\end{itemize}

\section*{Current Linking Rule}

The current best rule is:

\begin{quote}
\code{M6\_high\_precision\_gated\_score}
\end{quote}

It computes:

\[
\text{score} =
100 \times (
0.50 \times \text{buyer}
+ 0.25 \times \text{text}
+ 0.15 \times \text{CPV}
+ 0.05 \times \text{time}
+ 0.05 \times \text{geo}
)
\]

The selected setting is:

\begin{center}
\small
\begin{tabularx}{0.82\textwidth}{Xr}
\toprule
Parameter & Value \\
\midrule
Weight preset & \code{buyer\_precision\_v1} \\
Buyer weight & 0.50 \\
Text weight & 0.25 \\
CPV/theme weight & 0.15 \\
Time plausibility weight & 0.05 \\
Geography weight & 0.05 \\
Automatic threshold & 65 \\
Minimum text similarity & 0.25 \\
Text override for missing CPV & 0.55 \\
Minimum fuzzy-buyer text similarity & 0.45 \\
Minimum gap & 180 days \\
Maximum gap & 2,920 days \\
\bottomrule
\end{tabularx}
\end{center}

Acceptance requires:

\begin{itemize}
  \item strong buyer evidence: SIREN, normalized buyer name, or high fuzzy similarity;
  \item plausible future time gap;
  \item text similarity at least 0.25;
  \item CPV/theme continuity, unless text similarity is very high;
  \item score at least 65.
\end{itemize}

\section*{How the Weights Were Determined}

The current weights were not learned automatically by a statistical model. They were determined by:

\begin{enumerate}
  \item \textbf{Domain-informed design.} Buyer continuity should matter most for procurement renewal. Text similarity should matter strongly because the object of the contract identifies the need. CPV/theme should matter, but CPV can be noisy. Time matters, but duration is often missing. Geography matters least because the data is already restricted to Grand Ouest.
  \item \textbf{Pilot benchmark selection.} Several weight presets and thresholds were tested only on \code{PILOT\_DEVELOPMENT}. The selected setting was then evaluated on \code{LOCKED\_TEST}.
\end{enumerate}

Therefore the correct description is:

\begin{quote}
The linkage score is a transparent rule-based scoring model. Weights were specified from procurement-domain reasoning and selected through pilot-set tuning, then evaluated on a locked test subset.
\end{quote}

With more manually labelled pairs, future versions could estimate weights statistically using logistic regression, a Fellegi-Sunter model, gradient boosting, or active learning. With the current benchmark size, a transparent rule is more defensible than a complex model.

\section*{Parameter Diagnostics}

A dedicated parameter-selection notebook was created:

\begin{quote}
\pathcode{notebooks/09_linkage_parameter_selection_diagnostics.ipynb}
\end{quote}

It generated:

\begin{itemize}
  \item \pathcode{data/processed/boamp_grand_ouest/parameter_selection_grid.csv}
  \item \pathcode{data/processed/boamp_grand_ouest/parameter_selection_recommended_config.json}
  \item six figures under \pathcode{data/processed/boamp_grand_ouest/figures/parameter_selection/}
\end{itemize}

The current setting and diagnostic recommendation had the same pilot metrics:

\begin{center}
\small
\begin{tabularx}{\textwidth}{Xrr}
\toprule
Pilot metric & Current setting & Grid recommendation \\
\midrule
Coverage rate & 0.3125 & 0.3125 \\
Precision@1 & 0.800 & 0.800 \\
Recall@1 & 0.800 & 0.800 \\
Recall@5 & 1.000 & 1.000 \\
False-positive rate & 0.000 & 0.000 \\
No-link accuracy & 1.000 & 1.000 \\
Predicted count & 5 & 5 \\
\bottomrule
\end{tabularx}
\end{center}

The grid recommendation relaxed minimum text similarity from 0.25 to 0.20 on the pilot split, but because the pilot set contains only 16 anchors, the safer practical choice is to keep the current 0.25 threshold unless manual inspection supports relaxing it.

\section*{Why Precision Is Still a Concern}

The locked-test high-precision \metric{precision@1} is 0.4375. This means that among anchors where the model predicts a top successor, about 44\% of top predictions are correct in the locked benchmark.

That is not good enough to blindly treat all predicted links as true renewals. However, it is still useful if the project is framed correctly:

\begin{quote}
The method identifies a conservative subset of high-confidence observable renewals. It is not an exhaustive detector of all true renewals.
\end{quote}

For survival analysis, false positives are dangerous because a wrong renewal link creates a false event time. Therefore, the project should prioritize precision over coverage.

\section*{Current Limitations}

\begin{itemize}
  \item The benchmark is small: 94 eligible anchors, with only 16 in the pilot split.
  \item Current candidate generation is benchmark-anchor focused, not yet full-corpus episode linkage.
  \item BOAMP notices are not the same as contract episodes; initial notices, awards, modifications, lots, and linked notices may fragment one procurement need.
  \item Buyer identifiers are incomplete; most buyer keys are normalized names rather than SIREN.
  \item CPV can be noisy and is not sufficient alone for renewal detection.
  \item Duration and amount fields are often missing, limiting time-plausibility evidence.
  \item Survival analysis is not ready until the final renewal-link table is created.
\end{itemize}

\section*{Recommended Optimal Next Pipeline}

The most efficient next step is not to add many more algorithms. The next implementation should move toward the real research objective:

\begin{enumerate}
  \item Build procurement episodes from engineered notices.
  \item Restrict to Grand Ouest digital episodes.
  \item Generate candidate successor episodes for the full digital population.
  \item Apply the frozen high-precision linkage rule.
  \item Create \code{final\_renewal\_links.parquet}.
  \item Create \code{survival\_episodes.parquet}.
  \item Run survival analysis on high-confidence identifiable renewals.
  \item Add sensitivity analysis using stricter and looser linkage thresholds.
\end{enumerate}

The most important methodological improvement is:

\begin{quote}
Do not link raw notices directly if possible. First construct procurement episodes, then link episodes.
\end{quote}

\section*{Future Survival Metrics}

After linkage is frozen, survival analysis should use:

\begin{itemize}
  \item renewal observed indicator;
  \item time to renewal in months;
  \item right-censoring indicator;
  \item renewal probability within 12, 24, and 36 months;
  \item Kaplan-Meier survival curves by technology segment;
  \item log-rank tests across segments;
  \item Cox proportional hazards model coefficients and hazard ratios;
  \item sensitivity analysis across linkage thresholds.
\end{itemize}

\section*{Important Local Artifacts}

\noindent Most processed artifacts are under:

\begin{quote}
\pathcode{data/processed/boamp_grand_ouest/}
\end{quote}

\small
\begin{longtable}{p{0.34\textwidth}p{0.56\textwidth}}
\toprule
Purpose & File name \\
\midrule
\endhead
Engineered notices & \pathcode{notices_engineered.parquet} \\
Reference anchors & \pathcode{reference_anchor_episodes.parquet} \\
Reference links & \pathcode{reference_successor_links.parquet} \\
Candidate pairs & \pathcode{renewal_candidate_pairs.parquet} \\
Baseline predictions & \pathcode{linkage_baseline_predictions.parquet} \\
Baseline evaluation & \pathcode{linkage_evaluation_results.csv} \\
High-precision predictions & \pathcode{high_precision_linkage_predictions.parquet} \\
High-precision evaluation & \pathcode{high_precision_linkage_evaluation.csv} \\
High-precision config & \pathcode{high_precision_linkage_config.json} \\
Parameter grid & \pathcode{parameter_selection_grid.csv} \\
Parameter recommendation & \pathcode{parameter_selection_recommended_config.json} \\
Manual review candidates & \pathcode{high_precision_review_candidates.csv} \\
\bottomrule
\end{longtable}
\normalsize

\section*{Figure Folders}

The session produced three main diagnostic figure folders:

\begin{itemize}
  \item \pathcode{data/processed/boamp_grand_ouest/figures/linkage_baselines/}
  \item \pathcode{data/processed/boamp_grand_ouest/figures/high_precision_linkage/}
  \item \pathcode{data/processed/boamp_grand_ouest/figures/parameter_selection/}
\end{itemize}

The parameter-selection figures are:

\begin{itemize}
  \item \pathcode{01_pilot_threshold_metric_curves.png}
  \item \pathcode{02_pilot_precision_coverage_frontier.png}
  \item \pathcode{03_pilot_precision_heatmap_threshold_text.png}
  \item \pathcode{04_pilot_false_positive_heatmap_threshold_text.png}
  \item \pathcode{05_pilot_recall5_heatmap_threshold_text.png}
  \item \pathcode{06_pilot_current_vs_recommended_metrics.png}
\end{itemize}

\section*{Recommended Report Framing}

For the internship report, use this framing:

\begin{quote}
Because BOAMP does not provide explicit renewal labels, the study uses a conservative record-linkage approach. Candidate successor contracts are generated through temporal, buyer, CPV/theme, and text-based blocking. A transparent high-precision scoring rule is tuned on a pilot manual benchmark and evaluated on a locked test benchmark. The resulting links are interpreted as identifiable high-confidence renewals rather than exhaustive ground truth. Survival analysis is then performed on this conservative linked episode dataset, with censoring and sensitivity analysis used to account for unobserved or missed renewals.
\end{quote}

\section*{Chronological Session Log}

\begin{enumerate}
  \item Built raw BOAMP acquisition logic and kept 2026 separate from the 2015--2025 historical corpus.
  \item Created a full raw CSV and then a Grand Ouest raw CSV.
  \item Created annual Grand Ouest Parquet files.
  \item Built the EDA notebook and added a Grand Ouest section.
  \item Read the internship guide and reframed the project around renewal detection and survival analysis.
  \item Discussed why BOAMP has no ground truth and why manual benchmark evaluation is needed.
  \item Created preprocessing and data-engineering notebooks.
  \item Prepared the manual reference benchmark from three CSV files.
  \item Generated renewal candidate pairs for benchmark anchors.
  \item Compared baseline linkage algorithms.
  \item Diagnosed baseline false positives and false negatives.
  \item Implemented a stricter high-precision linkage iteration.
  \item Added parameter-selection diagnostic plots.
  \item Clarified that the project is currently at the linkage-method selection stage, not preprocessing or survival analysis.
  \item Reframed the next optimal pipeline around episode construction, conservative linkage, and survival analysis.
\end{enumerate}

\section*{Final Recommendation}

The next implementation should be:

\begin{quote}
Build procurement episodes from the engineered Grand Ouest digital notices, apply the frozen high-precision linkage rule to episode-level candidates, create a conservative renewal-link table, and then build the survival dataset.
\end{quote}

The current linkage method is useful, but it should be described honestly as a conservative high-confidence renewal identification method, not a complete renewal detector.

\end{document}
